\begin{lexlean}{Softplus}
\useglossary{lexlean.std.bool@1.2.0}
\useglossary{lexlean.std.nat@1.2.0}
\importmodule{ScalarF32}
\title{Natural number}
\begin{semanticmodule}
inductive{
  name: Request,
  type_parameters: [],
  parameters: [],
  constructors: [
    {name: input, fields: [Nat]},
    {name: exponential, fields: [Nat]},
    {name: logarithm, fields: [Nat, Nat, Nat]}
  ]
}

inductive{
  name: Step,
  type_parameters: [],
  parameters: [],
  constructors: [
    {name: done, fields: [Nat]},
    {name: exponentiate, fields: [Nat]},
    {name: takeLog, fields: [Nat, Nat]}
  ]
}

definition{
  name: evaluate,
  parameters: [request : Request],
  result: Step,
  body: (match request
    | Request.input word => (if ((1101004800 <? word) && (word <=? 2139095040)) then (ctor Step.done word) else (ctor Step.exponentiate word))
    | Request.exponential word => (let unit : Nat := (ScalarF32::add 1065353216 word); (if (unit =? 1065353216) then (ctor Step.done word) else (ctor Step.takeLog unit word)))
    | Request.logarithm unit value log => (let increment : Nat := (ScalarF32::add unit 3212836864); (let residual : Nat := (ScalarF32::add value (ScalarF32::negate increment)); (ctor Step.done (ScalarF32::add log (ScalarF32::divide residual unit)))))),
  executable: true,
  production: {effects: [overflow], targets: [rust-core]}
}

theorem{
  name: tinyExponentialSurvivesRoundedUnit,
  parameters: [],
  statement: ((evaluate (ctor Request.exponential 1)) = (ctor Step.done 1)),
  proof: reflexivity{}
}

theorem{
  name: aboveCutoffKeepsTheInput,
  parameters: [],
  statement: ((evaluate (ctor Request.input 1101004801)) = (ctor Step.done 1101004801)),
  proof: reflexivity{}
}

theorem{
  name: cutoffStillUsesExponential,
  parameters: [],
  statement: ((evaluate (ctor Request.input 1101004800)) = (ctor Step.exponentiate 1101004800)),
  proof: reflexivity{}
}
\end{semanticmodule}
\end{lexlean}
